% =====================================================================
% Section 3 — Design and Architecture
% =====================================================================
\section{Design and Architecture}\label{sec:design}

\subsection{Principles}

Three principles govern the design of JxMVC and explain most of its
concrete decisions.

\begin{description}
  \item[Zero runtime dependencies.] Every capability is implemented on
    top of the JDK and the Jakarta Servlet API. No third-party library
    is required at runtime; the only additional production artifact is
    the JDBC driver of the target database. This bounds the artifact
    size (a \SI{253}{\kilo\byte} JAR of 54 classes) and eliminates the
    transitive-dependency supply-chain surface.
  \item[Security-by-default.] Protections are properties of the request
    pipeline, not of individual routes. A newly written controller is
    hardened before its author adds a single security annotation;
    relaxing a protection is explicit and local.
  \item[Explicit, readable API.] The framework favours direct,
    inspectable control flow over ``magic''. Data access reads columns
    directly from a result row rather than mapping to reflection-bound
    POJOs, and the entire core is small enough to be read and audited
    end to end---a property we treat as a first-class feature for
    teaching and for security review.
\end{description}

\subsection{The request pipeline}
\label{sec:pipeline}

The heart of JxMVC is a fixed, ordered request pipeline implemented in
the core servlet (\texttt{MainLxServlet}). Every request traverses the
same fifteen-stage sequence, shown in its real execution order in
Figure~\ref{fig:pipeline}, so cross-cutting concerns---sanitisation,
CORS, rate limiting, authentication, validation, compression,
metrics---are applied uniformly and by construction. Grouping by role,
the principal stages are:

% Figura del pipeline de 15 etapas (orden real de MainLxServlet.serviceInternal).
% Requiere: \usepackage{tikz}\usetikzlibrary{positioning,arrows.meta} (en main.tex).
\begin{figure}[tb]
  \centering
  \small
  \begin{tikzpicture}[
    node distance=1.9mm,
    stage/.style={draw, rounded corners=1.5pt, minimum width=62mm,
      minimum height=4.9mm, inner sep=1pt, align=center, font=\footnotesize, fill=black!3},
    sec/.style={stage, fill=red!7, draw=red!55, very thick},
    arr/.style={-{Latex[length=1.4mm]}, gray!70}
  ]
    \node[stage] (s1)  {\textbf{1}\quad Internal endpoints \texttt{/jx/*}};
    \node[stage, below=of s1]  (s2)  {\textbf{2}\quad Metrics timer start};
    \node[stage, below=of s2]  (s3)  {\textbf{3}\quad Route resolution (\texttt{BaseDispatcher})};
    \node[sec,   below=of s3]  (s4)  {\textbf{4}\quad Rate limiting (per resolved action)};
    \node[sec,   below=of s4]  (s5)  {\textbf{5}\quad Authentication (fail-closed)};
    \node[sec,   below=of s5]  (s6)  {\textbf{6}\quad CSRF token (mutating verbs)};
    \node[sec,   below=of s6]  (s7)  {\textbf{7}\quad CORS policy};
    \node[stage, below=of s7]  (s8)  {\textbf{8}\quad Build context (\texttt{JxRequest}/\texttt{JxResponse})};
    \node[stage, below=of s8]  (s9)  {\textbf{9}\quad \texttt{before} filters};
    \node[stage, below=of s9]  (s10) {\textbf{10}\quad Controller instantiation + DI};
    \node[stage, below=of s10] (s11) {\textbf{11}\quad \texttt{@JxBeforeAction}};
    \node[stage, below=of s11] (s12) {\textbf{12}\quad \texttt{@JxModelAttr}};
    \node[stage, below=of s12] (s13) {\textbf{13}\quad Invoke (async/retry, cache, tx)};
    \node[stage, below=of s13] (s14) {\textbf{14}\quad \texttt{@JxAfterAction} + \texttt{after} filters};
    \node[stage, below=of s14] (s15) {\textbf{15}\quad Render + metrics};

    \foreach \a/\b in {s1/s2,s2/s3,s3/s4,s4/s5,s5/s6,s6/s7,s7/s8,s8/s9,
                       s9/s10,s10/s11,s11/s12,s12/s13,s13/s14,s14/s15}
      \draw[arr] (\a) -- (\b);

    \node[above=1.2mm of s1, font=\footnotesize\itshape] {request};
    \node[below=1.2mm of s15, font=\footnotesize\itshape] {response};
  \end{tikzpicture}
  \caption{The fixed fifteen-stage request pipeline of
    \texttt{MainLxServlet} (real execution order). Cross-cutting to every
    response, and applied outside this sequence, are the security headers
    and conditional gzip set in \texttt{service}. The shaded stages
    (4--7) form the security gate that precedes and encloses the
    developer's controller, which is what makes the ``secure-by-default''
    property structural rather than advisory.}
  \label{fig:pipeline}
\end{figure}

\begin{enumerate}
  \item \textbf{Sanitiser} (\texttt{BaseSanitizer}) --- XSS-oriented
    cleaning of request parameters and arguments.
  \item \textbf{CORS resolver} (\texttt{BaseCorsResolver}) ---
    declarative, annotation-driven CORS policy.
  \item \textbf{Dispatcher} (\texttt{BaseDispatcher}) --- routing by
    annotations, path templates \texttt{\{var\}}, and the
    \texttt{/controller/action/args} convention, with correct
    \texttt{405} handling and \texttt{HEAD}$\rightarrow$\texttt{GET}.
  \item \textbf{Rate limiter} (\texttt{JxRateLimiter}) --- per-client
    rate limiting keyed by action.
  \item \textbf{Security} (\texttt{JxSecurity}) --- authentication and
    role checks via a pluggable \texttt{JxAuthProvider}, fail-closed
    when no provider is registered.
  \item \textbf{Filters} (\texttt{JxFilters}) --- global
    before/after hooks.
  \item \textbf{Controller} --- the developer's action, with access to
    data via \texttt{JxDB}, \texttt{JxRepository}, and
    \texttt{JxTransaction}.
  \item \textbf{Validation} (\texttt{JxValidation}) --- constraint
    checking of bound request bodies.
  \item \textbf{Result} (\texttt{ActionResult}) --- a
    \textsc{view}/\textsc{text}/\textsc{json}/\textsc{redirect}
    outcome.
  \item \textbf{Compression} (\texttt{JxGzip}) --- gzip with a
    threshold-based passthrough.
  \item \textbf{Metrics} (\texttt{JxMetrics}) --- per-route counters
    and timing.
\end{enumerate}

Because the security-relevant stages (sanitisation, rate limiting,
authentication, validation, body caps) precede and enclose the
developer's controller, the ``secure-by-default'' property is
structural rather than advisory.

\subsection{Subsystems}

% Figura de arquitectura por capas. Requiere tikz + positioning,arrows.meta,fit (main.tex).
\begin{figure}[tb]
  \centering
  \resizebox{\columnwidth}{!}{%
  \begin{tikzpicture}[
    font=\footnotesize,
    col/.style={draw=black!35, rounded corners=2pt, align=left, inner sep=2mm,
      text width=40mm, minimum height=34mm, fill=white},
    band/.style={draw=black!40, rounded corners=2pt, align=center, inner sep=2mm,
      text width=176mm},
    arr/.style={-{Latex[length=1.6mm]}, gray!60}
  ]
    \node[col] (c1) {\textbf{HTTP \& MVC}\\[2pt]
      \texttt{MainLxServlet}\\ \texttt{BaseDispatcher}\\ \texttt{JxController}\\
      \texttt{JxRequest}/\texttt{JxResponse}\\ \texttt{ActionResult}\\ \texttt{JxFilters}, \texttt{JxGzip}};
    \node[col, right=3mm of c1] (c2) {\textbf{Data access}\\[2pt]
      \texttt{JxDB}\\ \texttt{JxPool}\\ \texttt{JxRepository}\\ \texttt{JxTransaction}\\
      \texttt{DBRow}/\texttt{DBRowSet}};
    \node[col, right=3mm of c2] (c3) {\textbf{Infrastructure}\\[2pt]
      \texttt{JxJson}, \texttt{GenApi}\\ \texttt{JxCache}\\ \texttt{JxScheduler}\\
      \texttt{JxEventBus}\\ \texttt{JxMetrics}, \texttt{JxOpenApi}\\ \texttt{JxWebSocket}, \texttt{JxValidation}};
    \node[col, right=3mm of c3, fill=red!5, draw=red!45] (c4) {\textbf{Security (by default)}\\[2pt]
      \texttt{JxSecurity}\\ \texttt{JxOAuth} (OIDC)\\ \texttt{JxPasswords}\\ \texttt{JxCsrf}\\
      \texttt{JxRateLimiter}\\ \texttt{JxHtml}, \texttt{BaseSanitizer}};

    \node[draw=none, inner sep=0, fit=(c1)(c4)] (cols) {};
    \node[band, fill=blue!7, above=3mm of cols] (app)
      {\textbf{Application} --- controllers (\texttt{@JxControllerMapping}), JSP views, DB-backed models};
    \node[band, fill=black!7, below=3mm of cols] (found)
      {\textbf{Foundation} --- JDK 17+ (\texttt{java.net.http}, \texttt{java.security}, reflection, Virtual Threads)
       \; + \; Jakarta Servlet / JSP / WebSocket API \textit{(provided scope)}};

    \draw[arr] (app.south) -- (app.south |- cols.north);
    \draw[arr] (found.north |- cols.south) -- (found.north);
  \end{tikzpicture}}
  \caption{Layered architecture of JxMVC. Every subsystem---HTTP/MVC,
    data access, infrastructure, and security---is implemented in-framework
    on top of the JDK and the \texttt{provided}-scope Jakarta Servlet API,
    with \textbf{zero third-party runtime dependencies} in the
    \SI{253}{\kilo\byte} core. The security column is not a bolt-on module:
    its checks are wired structurally into the request pipeline of
    Figure~\ref{fig:pipeline}.}
  \label{fig:architecture}
\end{figure}

Figure~\ref{fig:architecture} shows how JxMVC layers its subsystems on
the JDK and the Jakarta Servlet API. All subsystems are built
in-framework. Data access (\texttt{JxDB})
issues JDBC directly with named parameters and returns row objects
whose columns are read explicitly, avoiding reflection-based mapping.
\texttt{JxPool} is a custom connection pool with configurable
revalidation; \texttt{JxRepository} adds generic CRUD with soft delete
and pagination; \texttt{JxTransaction} provides thread-local JDBC
transactions. \texttt{JxJson} is a from-scratch JSON parser and
serialiser with input limits, complemented by a variadic JSON builder
(\texttt{GenApi}). \texttt{JxCache} is an in-memory cache with TTL and
an entry cap; \texttt{JxScheduler} runs cron, fixed-rate, fixed-delay,
and one-shot tasks; \texttt{JxEventBus} is a synchronous event bus.
\texttt{JxMetrics} exposes per-route statistics, \texttt{JxOpenApi}
generates an OpenAPI~3.0 specification automatically, and
\texttt{JxWebSocket} provides WebSocket endpoints with rooms and
broadcast. On Java~21 and later, the framework detects virtual threads
for asynchronous work.

\subsection{Security by default}

Beyond the pipeline ordering, v3.4.0 ships a set of hardening measures
that are on (or fail-closed) by default and tunable through
configuration properties: security headers, fail-closed authentication
in the absence of a provider, per-action rate limiting (so an attacker
cannot evade limits by rotating a path variable), session-token CSRF
protection, request-body and JSON size caps that return \texttt{413}
when exceeded, anti open-redirect enforcement, HTML output encoding for
XSS defence at render time, and PBKDF2-SHA256 password hashing with
constant-time verification. Internal introspection endpoints
(\texttt{/jx/info}, \texttt{/jx/metrics}, \texttt{/jx/openapi}) return
\texttt{404} under the production profile unless explicitly exposed,
and the health endpoint omits sensitive fields. The implementation of
the most security-critical of these---rate limiting and the OIDC
client---is described in Section~\ref{sec:impl}.
